% GENERATED FILE. Edit canonical lessons, then run npm run generate:books.
% canonical-source-hash: fnv1a64:31908bf8e2746f85
% canonical-lessons: SA-C111-vatya, SA-C111-niharah, SA-C111-sakha, SA-C111-parnam, SA-C111-kusumam

\chapter{Storm, Mist, Branch, Leaf, Flower}
\label{ch:sa-storm-mist-branch-leaf-flower}
\hlchaptermodality{111}

\begin{chapteropening}
\textbf{By the end of this chapter:} \emph{I can say a storm, mist, a branch, a leaf and a flower.}

Complete the last lesson of chapter 111: i can say a storm, mist, a branch, a leaf and a flower.
\end{chapteropening}

\section[\texorpdfstring{\sk{वात्या} (vātyā)}{वात्या (vātyā)}]{\sk{वात्या} (vātyā) — a storm}

\label{lesson:SA-C111-vatya}



\begin{quote}
Before the new one: say the Sanskrit for lightning, then the Sanskrit for thunder.
\end{quote}

\subsection*{You'll want to know: \sk{वात्या}}
\textbf{\sk{वात्या}} — \emph{vātyā} — \textquotedblleft{}a storm\textquotedblright{}.

\begin{grammarlens}[title={What you've built}]
One more word for this chapter.
\end{grammarlens}

\subsection*{Your turn}
\begin{itemize}
\raggedright
  \item \emph{Say it:} \emph{vātyā}
  \item \emph{Say it:} \emph{vātyā}, once more
  \item \emph{Say it:} say \emph{vātyā}
  \item \emph{Recall:} say the Sanskrit for lightning, then the Sanskrit for thunder, then say \emph{vātyā} again
\end{itemize}

\begin{tcolorbox}[breakable,colback=teal!4,colframe=teal!35!black,title={Before you move on}]
What does \sk{वात्या} mean? (\textquotedblleft{}A storm\textquotedblright{}.) Say it once more.
\end{tcolorbox}
\section[\texorpdfstring{\sk{नीहारः} (nīhāraḥ)}{नीहारः (nīhāraḥ)}]{\sk{नीहारः} (nīhāraḥ) — mist, fog}

\label{lesson:SA-C111-niharah}



\begin{quote}
Before the new one: say the Sanskrit for thunder, then the Sanskrit for a storm.
\end{quote}

\subsection*{You'll want to know: \sk{नीहारः}}
\textbf{\sk{नीहारः}} — \emph{nīhāraḥ} — \textquotedblleft{}mist, fog\textquotedblright{}.

\begin{grammarlens}[title={What you've built}]
One more word for this chapter.
\end{grammarlens}

\subsection*{Your turn}
\begin{itemize}
\raggedright
  \item \emph{Say it:} \emph{nīhāraḥ}
  \item \emph{Say it:} \emph{nīhāraḥ}, once more
  \item \emph{Say it:} say \emph{nīhāraḥ}
  \item \emph{Recall:} say the Sanskrit for thunder, then the Sanskrit for a storm, then say \emph{nīhāraḥ} again
\end{itemize}

\begin{tcolorbox}[breakable,colback=teal!4,colframe=teal!35!black,title={Before you move on}]
What does \sk{नीहारः} mean? (\textquotedblleft{}Mist, fog\textquotedblright{}.) Say it once more.
\end{tcolorbox}
\section[\texorpdfstring{\sk{शाखा} (śākhā)}{शाखा (śākhā)}]{\sk{शाखा} (śākhā) — a branch}

\label{lesson:SA-C111-sakha}



\begin{quote}
Before the new one: say the Sanskrit for a storm, then the Sanskrit for mist.
\end{quote}

\subsection*{You'll want to know: \sk{शाखा}}
\textbf{\sk{शाखा}} — \emph{śākhā} — \textquotedblleft{}a branch\textquotedblright{}.

\begin{grammarlens}[title={What you've built}]
One more word for this chapter.
\end{grammarlens}

\subsection*{Your turn}
\begin{itemize}
\raggedright
  \item \emph{Say it:} \emph{śākhā}
  \item \emph{Say it:} \emph{śākhā}, once more
  \item \emph{Say it:} say \emph{śākhā}
  \item \emph{Recall:} say the Sanskrit for a storm, then the Sanskrit for mist, then say \emph{śākhā} again
\end{itemize}

\begin{tcolorbox}[breakable,colback=teal!4,colframe=teal!35!black,title={Before you move on}]
What does \sk{शाखा} mean? (\textquotedblleft{}A branch\textquotedblright{}.) Say it once more.
\end{tcolorbox}
\section[\texorpdfstring{\sk{पर्णम्} (parṇam)}{पर्णम् (parṇam)}]{\sk{पर्णम्} (parṇam) — a leaf}

\label{lesson:SA-C111-parnam}



\begin{quote}
Before the new one: say the Sanskrit for mist, then the Sanskrit for a branch.
\end{quote}

\subsection*{You'll want to know: \sk{पर्णम्}}
\textbf{\sk{पर्णम्}} — \emph{parṇam} — \textquotedblleft{}a leaf\textquotedblright{}.

\begin{grammarlens}[title={What you've built}]
One more word for this chapter.
\end{grammarlens}

\subsection*{Your turn}
\begin{itemize}
\raggedright
  \item \emph{Say it:} \emph{parṇam}
  \item \emph{Say it:} \emph{parṇam}, once more
  \item \emph{Say it:} say \emph{parṇam}
  \item \emph{Recall:} say the Sanskrit for mist, then the Sanskrit for a branch, then say \emph{parṇam} again
\end{itemize}

\begin{tcolorbox}[breakable,colback=teal!4,colframe=teal!35!black,title={Before you move on}]
What does \sk{पर्णम्} mean? (\textquotedblleft{}A leaf\textquotedblright{}.) Say it once more.
\end{tcolorbox}
\section[\texorpdfstring{\sk{कुसुमम्} (kusumam)}{कुसुमम् (kusumam)}]{\sk{कुसुमम्} (kusumam) — a flower}

\label{lesson:SA-C111-kusumam}



\begin{quote}
Before the new one: say the Sanskrit for a branch, then the Sanskrit for a leaf.
\end{quote}

\subsection*{You'll want to know: \sk{कुसुमम्}}
\textbf{\sk{कुसुमम्}} — \emph{kusumam} — \textquotedblleft{}a flower\textquotedblright{}.

\begin{grammarlens}[title={What you've built}]
One more word for this chapter.
\end{grammarlens}

\subsection*{Your turn}
\begin{itemize}
\raggedright
  \item \emph{Say it:} \emph{kusumam}
  \item \emph{Say it:} \emph{kusumam}, once more
  \item \emph{Say it:} say \emph{kusumam}
  \item \emph{Recall:} say the Sanskrit for a branch, then the Sanskrit for a leaf, then say \emph{kusumam} again
\end{itemize}

\begin{tcolorbox}[breakable,colback=teal!4,colframe=teal!35!black,title={Before you move on}]
What does \sk{कुसुमम्} mean? (\textquotedblleft{}A flower\textquotedblright{}.) Say it once more.
\end{tcolorbox}
